% engines: lualatex, pdflatex
% expect-warnings: 0
% expect-errors: 0
%
% goban/blank: the intersections that are not part of the board have no lines and no star points; a line
% joins only two neighbours of the board, so next to a blank point it ends as at the edge. The stones and
% the background do not change. With a grid image, the image is only in the cells of the board.
\documentclass{article}
\usepackage[margin=1cm]{geometry}
\usepackage{../../latex/pgf-go}
\graphicspath{{../../doc/img/}}
\newcommand\case[1]{\par\medskip\noindent\textbf{#1}\par\nopagebreak}
\pagestyle{empty}
\begin{document}

\case{blank = E: two boards side by side (the star point E5 goes too); then a rectangle and a row}
\begin{goban}[9, blank = E]
	\stones[b]{C3, D4, D5}
	\stones[w]{F3, F4, G5}
\end{goban}
\begin{goban}[9, blank = {A9:C7, 5}]
	\stones[b]{D8, D6}
	\stones[w]{C6}
\end{goban}

\case{Single points: a hole in the middle, and an intersection left alone (no line reaches it)}
\begin{goban}[9, blank = {E5, A2, B1}]
	\stones[b]{E4, E6}
\end{goban}

\case{With partial and labels}
\begin{goban}[9, blank = {D, 6}, partial = A1:G8, label]
	\stones[w]{C3, E3}
\end{goban}

\case{With a grid image (the image is cut to the cells of the board) and in \textbackslash pgfgoset}
\begin{goban}[9, blank = E, grid = {image = {file = pgfgo-grid.png}}]
	\stones[b]{C3}
\end{goban}
\pgfgoset{goban = {blank = {A, J}}}
\begin{goban}[9]
	\stones[b]{B2}
\end{goban}

\end{document}
